% A report fragment; requires amsmath, amssymb, and booktabs.
\section{Command and source audit for the two new composites}
\label{sec:composite-command-audit}

The two composites combine the supplied diagonal-vacuity axiom emitters
with the supplied changes that make Shift test the cached head alone and
make successful modus ponens return its result in that cache. These changes
are compatible. The ten-entry emitter requires changed stack indices;
the nine-entry command compiler cannot simply be copied verbatim.
This section proves the required command identities uniformly, binds the
changed primitive source tails, and states the limits of the audit.

\subsection{Exact command interface and emitted stack}

Let $p(a,b)=\binom{a+b+1}{2}+b$, let $[a;w]=p(a,w)$, and write
$[a_1,\ldots,a_k;w]$ for iterated cons with $a_1$ at the top.
The total Cantor inverse satisfies
\[
 \operatorname{hd}[a;w]=a,\qquad
 \operatorname{tl}[a;w]=w,
 \qquad \operatorname{hd}(0)=\operatorname{tl}(0)=0.
\]
In particular, $[0;0]=0$. No argument below assumes that a zero push
increases a stored stack length, or that an empty stack can be distinguished
from a stack of zero entries. The identities are numerical identities for
the total projections, valid also at this alias.

Use $I(a,b)=4p(a,b)$ and $U(a,B)=4p(B,a)+2$, and write a command store
as $\langle t,w,d,h\rangle$. The two modified commands are
\[
 S\langle t,w,d,h\rangle=
 \begin{cases}
  \mathsf{halt},&h=1,\\
  \langle h,\operatorname{tl}(w),d,\operatorname{hd}(w)\rangle,&h\ne1,
 \end{cases}
\]
and, putting $y=d+x$,
\[
 \mathsf{MP}_x\langle t,w,d,h\rangle=
 \begin{cases}
  \langle0,w,0,y\rangle,&I(t,y)=h,\\
  \langle0,w,y,0\rangle,&I(t,y)\ne h.
 \end{cases}
\]
Only $S$ can halt. Neither branch of MP pushes a stack entry.

Let $n$ be the \emph{total} number of entries AX pushes, including the
saved top: $n=9$ for the diagonal composite and $n=10$ for the separate
diagonal-head composite. For $D=d+x$, $H=h+y$, $A=c$, the effect is
\[
 \mathsf{AX}_{c,x,y}\langle t,w,d,h\rangle
 =\langle K(H,D,c),[q_1,\ldots,q_{n-1},t;w],0,0\rangle.
\]
Writing $L_1,\ldots,L_6$ for the original six heads in their push order,
the generated entries, listed from the stack top downward, are exactly
\[
\begin{array}{c|l}
n& (q_1,\ldots,q_{n-1})\\ \hline
9&(L_6,L_5,L_4,L_3,L_2,L_1,\rho_8(D,A,H,d),U(d,h))\\
10&(\Gamma(A),L_6,L_5,L_4,L_3,L_2,L_1,
                 \rho_8(D,A,H,d),U(d,h)).
\end{array}
\]
Here $\Gamma(B)=I(B,U(B,B))$. In the first case $K$ is the guarded
bundle $F_\Gamma=\neg(\Gamma(c)\to(V\to Q))$; in the second it is
the guarded old gate $G=\neg(V\to Q)$. A failed numerical guard gives
$K=0$, while the separately pushed $\Gamma(A)$ remains present in the
second case. The emitter audit checks these exact free-algebra trees and
the complete stack order, not a list of sampled formula codes.

All $q_j$ except the last have residue zero modulo four; the last is
$U(d,h)$ and has residue two. The gate, including a failed guard, has
residue zero. Consequently none of these generated values is code one.
This property is enough for all the termination claims below; it does not
assume that a theorem of the underlying calculus is numerically nonzero.

\subsection{Uniform stack-word theorem}

Commands execute from left to right. For every incoming store, including
one with $t=1$, the following are total, nonhalting command executions:
\begin{align*}
 \mathsf{AX}_{c,x,y};S^{n-1}:
 \langle t,w,d,h\rangle
 &\longmapsto
 \langle\rho_8(d+x,c,h+y,d),[t;w],0,U(d,h)\rangle,\\
 \mathsf{AX}_{c,x,y};S^n:
 \langle t,w,d,h\rangle
 &\longmapsto \langle U(d,h),w,0,t\rangle.
\end{align*}
Indeed, the first Shift tests the cleared cache zero; Shifts two through
$n$ test $q_1$ through $q_{n-1}$. None tests the saved top $t$. This is
a finite induction using the two Cantor projection identities, so it also
covers $t=w=0$ and all other zero aliases.

Define
\begin{align*}
 \mathsf{Push}_n&=\mathsf{AX}_{0,0,0};S^{n-1};\mathsf{MP}_0,\\
 \mathsf{Pop}_n&=\mathsf{AX}_{0,0,0};S^{n+1},\\
 \mathsf{Discard}_n&=\mathsf{MP}_1;\mathsf{AX}_{1,0,0};S^n;
                    \mathsf{MP}_0.
\end{align*}
Their complete contracts are
\begin{align*}
 \mathsf{Push}_n\langle t,w,d,h\rangle
   &=\langle0,[t;w],0,0\rangle,\\
 \mathsf{Pop}_n\langle t,[a;w],d,h\rangle
   &=\begin{cases}
       \mathsf{halt},&t=1,\\
       \langle t,w,0,a\rangle,&t\ne1,
     \end{cases}\\
 \mathsf{Discard}_n\langle t,w,0,0\rangle
   &=\langle0,w,0,0\rangle.
\end{align*}
Push's last MP compares an implication, of residue zero, with $U(d,h)$,
of residue two, and therefore takes its failure branch. Pop's last Shift
tests precisely the saved top; this is the only possible halt in that
word. For Discard, $I(t,1)>0$ makes the first MP fail, producing
$\langle0,w,1,0\rangle$. The AX identity gives
$\langle U(1,0),w,0,0\rangle$, and $I(U(1,0),0)>0$ makes the final
MP fail. These arguments cover arbitrary $t$, including zero and one.

Thus the concrete exponents are
\[
\begin{array}{c|ccc}
 &\mathsf{Push}:S^{n-1}&\mathsf{Pop}:S^{n+1}
 &\mathsf{Discard}:S^n\\ \hline
\text{diagonal}&8&10&9\\
\text{separate diagonal head}&9&11&10.
\end{array}
\]
In particular, Pop is \emph{not} claimed total for arbitrary preserved
top: it deliberately detects a derived contradiction.

\subsection{Selectors and generalization}

Let $\mathsf{Offset}_0$ be empty and let
$\mathsf{Offset}_u=\mathsf{MP}_u$ for $u>0$. Then
\[
 \mathsf{Offset}_u\langle0,w,0,0\rangle=\langle0,w,u,0\rangle,
\]
since $I(0,u)>0$ when $u>0$. Define
\[
 \mathsf{Gen}_{i,n}=\mathsf{Push}_n;\mathsf{Offset}_i;S;
                  \mathsf{AX}_{i,0,0};S^n.
\]
For every $a,w$, including $a=1$,
\[
 \mathsf{Gen}_{i,n}\langle a,w,0,0\rangle
       =\langle U(i,a),w,0,0\rangle.
\]
The Shift before AX tests cache zero and retrieves $a$ into $h$; AX
clears that cache before any subsequent Shift and constructs $U(i,a)$.
The last identity from the preceding subsection then proves the result.

For $1\le j\le n-3$, the ordinary head selector is
\[
 \mathsf{Discard}_n;\mathsf{AX}_{A,D,H};S^{j+1};
                  \mathsf{Pop}_n^{\,n-1-j}.
\]
It returns $\langle q_j,w,0,0\rangle$ from
$\langle t,w,0,0\rangle$ for arbitrary $t,w$. After the displayed
Shifts, $t=q_j$, the cache contains $q_{j+1}$, and the stack is
$[q_{j+2},\ldots,q_{n-1},0;w]$. There are exactly $n-1-j$ entries
in that displayed stack segment. Each Pop preserves $q_j\ne1$, consumes
one entry and clears the previous cache before its own Shift tests.
The last consumed entry is the saved zero, giving cache zero.

For $n=10$, the new diagonal head is $j=1$, so the concrete selector
for $\Gamma(B)$ is
\[
 \mathsf{Discard}_{10};\mathsf{AX}_{B,0,0};S^2;
                           \mathsf{Pop}_{10}^{8}.
\]
This is valid even for a failed gate guard and even for $B=0$.
Each original head has index one larger than in the nine-entry layout.

For $u\le a$, the restricted propositional head is selected by
\[
 \mathsf{Discard}_n;\mathsf{Offset}_u;
 \mathsf{AX}_{b,a-u,c};S^{n-1};\mathsf{Pop}_n,
\]
which returns $\rho_8(a,b,c,u)$ with the original stack and zero caches.
The gate selector is
\[
 \mathsf{Discard}_n;\mathsf{AX}_{c,b,a};\mathsf{Pop}_n^{n}.
\]
It returns $K(a,b,c)$ with the original stack and zero caches. Every Pop
is safe because its preserved top is the gate, which has residue zero.
These selectors are total even when their incoming top is one.

\subsection{Modus ponens, early halting, and a uniform proof compiler}

The MP suffix is independent of the AX stack length:
\[
 \mathsf{MPSuffix}_b=S;S;\mathsf{MP}_b;S.
\]
For every $a,b,w$,
\[
 \mathsf{MPSuffix}_b
  \langle0,[a,I(a,b),0;w],0,0\rangle
 =\begin{cases}
   \mathsf{halt},&a=1\text{ or }b=1,\\
   \langle b,w,0,0\rangle,&a\ne1\text{ and }b\ne1.
  \end{cases}
\]
The first Shift retrieves $a$. The second either halts at $a=1$ or
sets the top to $a$ and the cache to $I(a,b)$. MP succeeds and puts $b$
in the cache. The last Shift either halts at $b=1$ or retrieves the dummy
zero, leaving the underlying stack unchanged. If $a=b=w=0$, all three
displayed entries alias the zero stack. The same projection calculation
still applies and MP succeeds at $I(0,0)=0$. Omitting the dummy zero would
consume the first underlying entry, so its presence is essential.

For a finite proof tree $p$ in the appropriate numerical calculus, compile
axioms by the selectors above and compile generalization by
$W_{p_0};\mathsf{Gen}_{i,n}$. For an MP node with minor premise $a$ and
major premise $I(a,b)$, use
\[
 W_p=\mathsf{Discard}_n;\mathsf{Push}_n;
 W_{p_1};\mathsf{Push}_n;
 W_{p_0};\mathsf{Push}_n;\mathsf{MPSuffix}_b.
\]
Induction on $p$ proves the exact invariant:
for every $t,w$, starting from $\langle t,w,0,0\rangle$, the compiled
word either halts or returns $\langle\operatorname{concl}(p),w,0,0\rangle$.
The first two words in the MP case introduce the dummy zero; the two
premise computations and their Pushes then establish the suffix's exact
three-entry precondition. Early halting is preserved by command-word
extension and causes no completeness failure.

If $p$ proves code one, the word
\[
 \mathsf{AX}_{0,0,0};W_p;\mathsf{Pop}_n
\]
halts from \emph{every} incoming store. The first AX clears $d,h$;
the compilation invariant tolerates its arbitrary resulting top and stack.
If no earlier command halts, the final Pop detects the proven top one
before inspecting the arbitrary underlying stack. This universal incoming
store property is exactly what is needed for the inherited exhaustive
command-word scheduler; earlier trials need not restore an empty store.

\subsection{Soundness invariant and scope of the calibration}

Maintain derivations in the chosen numerical calculus for $t$, $h$,
and each finitely decoded stack entry. No derivability requirement is
placed on the offset $d$. Code zero has the fixed derivation
$\rho_8(0,0,0,0)=0$, because $I(0,0)=0$. Thus the infinite zero tail
does not require infinitely much proof data: a single fixed zero proof
supplies every requested padding entry.

AX preserves this invariant: its saved top is already derived, its
ordinary heads and guarded gate are axioms or zero, the separate
$\Gamma(A)$ is an unconditional axiom, its propositional head satisfies
$d\le d+x$, and $U(d,h)$ follows by theorem generalization from the old
cache. A failed MP returns zero; a successful MP derives $d+x$ from the
derived $t$ and $I(t,d+x)=h$. Shift only moves derived formulas, and a
halting Shift exposes a derivation of one. Together with the compiler and
the inherited scheduler, this proves command-program halting if and only
if the selected numerical calculus has a refutation.

The compiler, selectors, finite proof-store updates, and the two changed
loop-trace constructors are primitive recursive. The correctness arguments
use ordinary induction on finite proofs, lists, natural loop counters and
traces. They therefore supply the changed-command part of a PA-internal
calibration. This audit does not independently reprove the entire inherited
PA-to-calculus and calculus-to-PA translation, nor claim a newly checked
Lean or Rocq endpoint. Those logical translations and the internal coding
of the literal tape simulation remain the inherited interfaces required
for an end-to-end assertion about $\operatorname{Con}(\mathrm{PA})$.

\subsection{Finite source binding and arbitrary-value loop proofs}

An independent pair of CFG resolvers checks the complete source graphs.
Let $N$ be the new program's number of primitive sites. Set
$E=N-15$, $F=N-12$, $S_0=N-11$. Its first $E$ primitive instructions
are identical to those of the corresponding old axiom layout, except for
three exit retargetings: source site $48$ branches to the new Shift entry,
the last AX site $N-96$ returns to $F$, and the last MP-prelude site
$N-16$ returns to $F$. No AX macro or formula-emission instruction changes.
The concrete entries are
\[
\begin{array}{lrrrr}
\toprule
 &N&\text{AX final}&\text{MP success }E&\text{Shift }S_0\\
\midrule
\text{diagonal composite}&3676&3580&3661&3665\\
\text{separate composite}&3649&3553&3634&3638\\
\bottomrule
\end{array}
\]
The new MP tail is the three-instruction loop
$\operatorname{dec}D;\operatorname{inc}H;\operatorname{jump}$ followed
by return to fetch site $25$. From $(D,H)=(y,0)$ it reaches $(0,y)$
in exactly $3y+1$ instructions and returns in one more. Failure returns
directly with $(D,H)=(y,0)$. This proves the displayed MP contract for
every natural $y$.

The inherited command-boundary convention gives scratch $X=0$;
the four stored values $t,w,d,h$ remain arbitrary. The new Shift first
clears $T$ in $2t+1$ instructions, then transfers
$H$ to $T$ and scratch $X$ in $4h+1$ instructions, then restores $H$
from $X$ in $3h+1$ instructions. Thus it reaches $(T,H,X)=(h,h,0)$
after exactly $2t+7h+3$ instructions. Its final two instructions set
$X=1$ and enter the unchanged guard, which compares $H$ to one and
otherwise calls the unchanged total Cantor decoder. The elementary loop
invariants prove these effects for every natural command-store value
under that boundary convention, not just for a bounded execution prefix.

\paragraph{Executable audit.}
The finite binding receipt is \texttt{composite\_binding.json}.
The command algebra supplement checks 164 symbolic cases for $n=9$ and
188 for $n=10$, with both opaque arbitrary tails and the exact zero-stack
alias; an undecided equality is rejected. Separately, the independent
whole-tape symbolic checker passes all 4,784 and 4,749 source outcomes and
the exact 57,922-step blank startup. These checks support the displayed
mathematical proofs and identify the literal implementations; they are
not substitutes for an end-to-end kernel-checked arithmetic theorem.
